\documentclass[11pt]{article}
\usepackage[OT1]{fontenc}
\renewcommand{\rmdefault}{cmr}
\renewcommand{\sfdefault}{cmss}
\renewcommand{\ttdefault}{cmtt}
\usepackage{amsmath,amssymb,amsthm,mathtools}
\usepackage[margin=27mm]{geometry}
\usepackage{microtype,booktabs,longtable}
\usepackage[hidelinks]{hyperref}
\setlength{\emergencystretch}{2em}
\setlength{\parskip}{3pt plus 1pt}
\setcounter{tocdepth}{1}
\newtheorem{theorem}{Theorem}[section]
\newtheorem{lemma}[theorem]{Lemma}
\newtheorem{proposition}[theorem]{Proposition}
\newtheorem{corollary}[theorem]{Corollary}
\theoremstyle{definition}
\newtheorem{definition}[theorem]{Definition}
\newtheorem{example}[theorem]{Example}
\newtheorem{problem}[theorem]{Open problem}
\theoremstyle{remark}
\newtheorem{remark}[theorem]{Remark}
\newcommand{\N}{\mathbb N}
\newcommand{\NN}{\mathbb N_{>0}}
\newcommand{\Ftwo}{\mathbb F_2}
\newcommand{\lv}{\operatorname{lv}}
\newcommand{\norm}{\operatorname{N}}
\newcommand{\sz}{\operatorname{sz}}
\newcommand{\can}[1]{\mathopen{[\![}#1\mathclose{]\!]}}
\newcommand{\mle}{\preccurlyeq}
\newcommand{\mslt}{\prec}
\newcommand{\PhiF}{\Phi}
\newcommand{\Esrc}{\mathcal E_{\mathrm{src}}}
\newcommand{\headpart}{\operatorname{lim}}
\newcommand{\wt}{\operatorname{wt}}
\newcommand{\Ty}{\operatorname{Ty}}
\newcommand{\FV}{\operatorname{FV}}
\newcommand{\NF}{\operatorname{NF}}
\newcommand{\sgn}{\operatorname{sign}}
\newcommand{\id}{\operatorname{id}}
\newcommand{\BH}{\mathsf{BH}}
\newcommand{\ord}{\operatorname{ord}}
\newcommand{\bodyurl}[1]{\url{#1}}
\numberwithin{equation}{section}

\newcommand{\ff}{f}
\newcommand{\JJ}{J}
\newcommand{\XX}{X}
\newcommand{\EE}{E}
\newcommand{\VV}{V}
\newcommand{\HH}{H}
\newcommand{\pp}{p}
\title{Multivariate Decidable Fragments\\of Canonical Comparison}
\author{Logic Lineages Working Papers}
\date{Working manuscript, October 2026}
\begin{document}
\maketitle
\begin{abstract}
The canonical ordinal assignment used by Wilken and Weiermann turns a
termination question for G\"odel's $T$ into universal comparison of numerical
collapse expressions. Several variables create a difficulty absent from
closed evaluation: derived values share inputs and cannot be replaced by
independent unknowns. We prove decision results which retain those
correlations. For the exact control $\Phi(n)=F_5(n+100)$, we verify the
hypotheses of Snyatkov's published arithmetic theorem. This settles a
finite-coefficient source class containing every recursor-free term with
ground free variables, regardless of its internal types. We then give exact
finite-degree descent kernels, ranked cancellation, and a terminating closed
existential procedure. These decide canonical comparison when every internal
type has level at most two. At arbitrary fixed ordinal heads, different
arguments give separate decision results for nested common-head sums,
balanced heterogeneous layers, and rows depending on one arbitrary positive
core. The last result reduces directly to a finite Presburger predicate.
All algorithms include their finite exceptions. None decides unrestricted
comparison of all ground-open canonical assignments.
\end{abstract}
\tableofcontents

\section{What is being compared?}
Wilken and Weiermann's assignment associates with a typed term a vector of
ordinal expressions. Their effectiveness question asks, among other things,
whether its universal comparison relations can be decided effectively
\cite{ww}. The assignment is computed from the original term, so two terms
with the same normal form can have different assigned vectors. We keep that
canonical assignment throughout; normalization is not a preliminary step.

This paper concerns terms whose free variables have ground type. For their
\emph{actual generated assignments}, the positive coordinates are closed
ordinals. The zeroth coordinate is a natural-valued expression made from
positive-natural ground variables, natural constants, addition, fixed natural
multiples, and functions
\[
 h_\lambda(n)=\psi(\lambda+n),
 \qquad 0<\lambda<\varepsilon_0\text{ a fixed limit}.
\]
The forward-translation theorem in the foundations companion
\cite{wwfoundations} proves this by simultaneous induction on generated
terms and abstraction operations. The qualifier ``actual generated''
matters: the unrestricted ambient classes of ordinal expressions in the
original definition also permit promotions of whole scalar expressions.
Those larger classes do not have the same closure invariant merely by
membership. We use the stronger proved invariant only for the source
assignments to which it applies.

Use the additive Cantor-normal-form norm
\[
 N(0)=0,\qquad
 N\!\left(\sum_i\omega^{\alpha_i}c_i\right)
       =\sum_i c_i\bigl(1+N(\alpha_i)\bigr).
\]
For a fixed limit set
\begin{equation}\label{eq:J}
 J_\lambda(t)=\psi\bigl(\lambda+t-N(\lambda)\bigr),
       \qquad t\ge N(\lambda),\qquad J_0(t)=t.
\end{equation}
Thus $t$ is the whole ordinal's norm, not just its finite tail. In particular,
an offset $N(\lambda)$ must not disappear when two heads are compared.
The foundations companion proves the exact closed evaluator, controlled
strictness, and the complete-count identity
\begin{equation}\label{eq:complete}
 J_\lambda(t)=Z+\tau,\qquad Z=f^\tau(t),\qquad \tau\ge1,
\end{equation}
where $f=\Phi$ and $\tau$ counts replacements until the ordinal becomes
finite. Finite successors after that point are not included in $\tau$.
These are the foundational facts imported here, with their proofs in the
included manuscript. The remaining comparison mechanisms are proved below.

Our results are complementary. A bound on internal type levels permits
arbitrary arrangements of polynomial heads. Common-head and balanced
certificates instead permit nonpolynomial heads but restrict how the open
arguments are arranged. Neither condition subsumes the other. Failure of a
certificate means that its theorem does not apply; it is never a negative
answer to the requested comparison.

\section{The exact control and a published arithmetic theorem}
Define, literally,
\begin{equation}\label{eq:f}
 F_0(n)=2^n,\qquad F_{j+1}(n)=F_j^{\,n+1}(n),\qquad
 f(n)=F_5(n+100).
\end{equation}
A superscript on a function denotes a finite iterate. The starting argument
in $F_{j+1}$ varies with $n$. This convention will not be replaced by
iteration from a fixed seed.

\begin{lemma}[Growth]\label{lem:growth}
Each $F_j$ is recursive and strictly increasing, with $F_j(n)\ge2^n$.
For $j\ge0$,
\[
 F_{j+1}(n+1)\ge2^{F_{j+1}(n)+1}.
\]
Consequently $f(z+1)\ge(A+2)f(z)$ whenever $z\ge A\ge0$.
\end{lemma}
\begin{proof}
Finite iteration preserves recursiveness, strict increase, and domination
of $2^n$. Increasing the input from $n$ to $n+1$ makes the first $n+1$
iterates at least one larger, and the definition then supplies one further
application of $F_j$. Thus
$F_{j+1}(n+1)\ge F_j(F_{j+1}(n)+1)\ge2^{F_{j+1}(n)+1}$.
For $v\ge1$, $2^{v+1}\ge(v+2)v$, by direct checking at $v=1,2,3$ and
induction. Apply this with $v=f(z)>z\ge A$.
\end{proof}

\begin{lemma}[Effective constant residues]\label{lem:residues}
For every $m\ge1$ one can compute $r_m\in[0,m)$ such that
\[
 F_j(n)\equiv r_m\pmod m\quad(j\ge1,\ n\ge m).
\]
In particular $f(n)\equiv r_m\pmod m$ for $n\ge\max(0,m-100)$.
For $s\ge1$, $j\ge2$, this can be strengthened to
\begin{equation}\label{eq:envelope}
 m\le 2^{2^s}\quad\Longrightarrow\quad F_j(s)\equiv r_m\pmod m.
\end{equation}
If $B=f(u)$ and $m\le2^{2^B}$, then also $f(u+1)\equiv r_m\pmod m$.
\end{lemma}
\begin{proof}
Set $r_1=0$. For $m>1$, write $m=2^a b$ with $b$ odd, and put
$q=\varphi(b)$ if $b>1$, $q=1$ otherwise. Since $q<m$, recursively
compute $r_q$, then use the Chinese remainder theorem to define
\[
 r_m\equiv0\pmod{2^a},\qquad r_m\equiv2^{r_q}\pmod b.
\]
The recursion strictly lowers the integer modulus.
For towers $T_0(n)=n$, $T_{h+1}(n)=2^{T_h(n)}$, induction on $m$ shows
$T_h(n)\equiv r_m\pmod m$ whenever $h,n\ge m$. Indeed, the last exponent
is at least $n\ge a$, while its residue modulo $q$ is $r_q$ by induction;
Euler's theorem handles the odd factor. This proves the assertion for
$F_1(n)=T_{n+1}(n)$. At every higher level, all intermediate arguments are
at least $n$, so the last application of the preceding level proves the
same residue assertion.

For the stronger bound, $F_1(s)\ge2^{2^s}$. In $F_2(s)$ there are at
least two $F_1$ applications; the final one has input at least $F_1(s)$,
so its input is at least $m$. The same argument propagates through every
$j\ge2$. Finally put $n=u+100$. The last $F_4$ call in
$F_5(n+1)=F_4^{n+2}(n+1)$ has input
$F_4^{n+1}(n+1)>F_4^{n+1}(n)=B$. Apply the stronger bound to that call.
\end{proof}

\begin{theorem}[Snyatkov's arithmetic supplier]\label{thm:snyatkov}
The first-order theory $\operatorname{Th}(\N;0,1,+,<,f)$, with the exact
function in \eqref{eq:f}, is decidable.
\end{theorem}
\begin{proof}
We use Theorem~2 and effective Corollary~2 of Snyatkov~\cite{snyatkov},
pp.~50--51. Its definitions on p.~40 require an increasing recursive
function effectively compatible with addition: effective eventual periodicity
modulo every integer, and an effective shifted sign separation for each
unbounded integer combination of fixed translates. It also requires effective
eventual periodicity of the \emph{fixed-seed} hyperfunction
$H(0)=1$, $H(n+1)=f(H(n))$. The cited theorem gives a decision procedure in
an expanded language with the stated auxiliary divisibility and inverse
symbols, hence in its $f$-only reduct.

The residue condition follows from Lemma~\ref{lem:residues}. For the second
condition combine equal shifts in $S(x)=\sum_b a_b f(x+b)$. If every
combined coefficient is zero the sum is zero. Otherwise let $B$ be the
largest shift with nonzero coefficient, and set
\[
 A=\sum_{b<B}|a_b|,\qquad
 \delta=\max(0,-\min b,A+1-B).
\]
Every argument at $x+\delta$ is nonnegative. With
$z=x+\delta+B-1\ge x+A$, Lemma~\ref{lem:growth} gives, if $a_B>0$,
\[
 S(x+\delta)\ge f(z+1)-Af(z)\ge2f(z)>f(x).
\]
For $a_B<0$, apply the same calculation to $-S$. This proves unboundedness,
its sign, and a suitable shift by finite arithmetic on the input data.
Lastly $H(n)\ge n+1$; for $n\ge m$, the last application of $f$ in
$H(n)$ has input at least $m$, so $H(n)\equiv r_m\pmod m$.
All published hypotheses are therefore paid. No identification of $H$ with
a diagonal $F_j$, and no decision result for variable-seed iteration, is
being imported.
\end{proof}

\section{An actual-source finite-coefficient class}
The foundations companion proves the identity
\begin{equation}\label{eq:finite-head}
 \psi(\omega k+n)=f^k(n+2k)+k\qquad(k,n\in\N).
\end{equation}
Although an iteration occurs here, its length $k$ is a closed integer
computed from source syntax. Replacing the expression by $k$ copies of $f$
is an effective finite translation.

\begin{theorem}[Finite-coefficient source comparison]\label{thm:finite-source}
Take two original well-typed $T$ terms of the same type with ground free
variables. Compute their canonical assignments. Suppose that at every
zeroth-coordinate collapse $\psi(\omega\otimes c+g)$ its closed coefficient
$c$ has finite ordinal value. Their universal weak, WW-strict, and equality
comparisons are decidable. Every recursor-free such term passes this test,
with no bound on its internal types.
\end{theorem}
\begin{proof}
Use the actual-source forward translation of \cite{wwfoundations};
$\otimes$ is natural product at this source interface. Evaluate closed
positive coordinates first. Weak vector comparison requires weakness in
all coordinates; WW-strict comparison requires strictness at coordinate
zero and weakness above it. Equality requires equality everywhere. A failed
positive-coordinate test already decides the requested universal relation.

For the remaining scalar expressions, retain each original variable name,
addition, and each fixed coefficient. Replace a collapse with finite
coefficient $k$ by $g\mapsto f^k(g+2k)+k$ using
\eqref{eq:finite-head}. This translation preserves its value at every
positive valuation. A ground vector has only one coordinate, so admissible
closed substitutions give exactly all positive integers; there is no hidden
coupling between distinct ground variables. Apply Theorem~\ref{thm:snyatkov}
to the required universal formula, with $x_i\ge1$ explicitly imposed.

The source theorem also proves that recursor-free actual assignments have
finite positive coefficients, by induction on their term and abstraction
construction. Thus all these sources pass the test. This argument uses that
actual-source theorem, not an assertion about all ambient ordinal-expression
classes. A recursor need not cause failure: the canonical vector of $R^oX$
is $(f(x+4)+1,1,\omega)$, whose direct scalar collapse coefficient is one.
Closed coefficient evaluation and the finite syntax traversal terminate.
\end{proof}

\section{Exact finite-degree prefix kernels}
The next class requires counters as well as terminal norms. A dominant
approximation would lose cancellations needed for multivariate comparison.
For integer $d\ge1$, retained norm $c\ge0$, and sufficiently large $t$, start
with suffix $\omega^d$. At step $i\ge1$ use total norm $f^i(t)$ and replace
the nonfinite suffix by the greatest smaller suffix of norm at most
$f^i(t)-c$. Stop as soon as it is finite. Let $E^d_c(t)$ be the number of
steps and
\[
 X^d_c(t)=f^{E^d_c(t)}(t).
\]
Each replacement strictly decreases an ordinal, so the process terminates.
The exact maximal-predecessor theorem in \cite{wwfoundations} supplies the
finite algorithm and full-budget saturation. The numerical definition makes
sense even for $c$ not realized by any actual ordinal prefix.

For a genuine prefix $\mu$ whose least CNF exponent is at least $d$, and
$c=N(\mu)$, splitting the controlled history gives
\begin{equation}\label{eq:prefix}
 J_{\mu+\omega^d}(t)=J_\mu(X^d_c(t))+E^d_c(t),
       \qquad t\ge c+d+1.
\end{equation}
The output of $X$ is still the \emph{total} norm, so the prefix mass is not
subtracted again. For fixed signed $j$, put
$X^d_{c,j}(t)=f^{E^d_c(t)+j}(t)$ on a tail where the exponent is positive.
These are coherent depths: $f(X^d_{c,j})=X^d_{c,j+1}$. A negative $j$ is a
fixed number of already available predecessors under $f$; it is not a new
variable-iteration graph.

\begin{lemma}[Exact block recursion]\label{lem:kernel-recursion}
One has
\[
 E^1_c=1,\quad X^1_c=f,\qquad
 E^2_c(t)=f(t)/2-\lceil c/2\rceil+1.
\]
Define
\begin{equation}\label{eq:V}
 V_j(t)=f^{\,f(t)/2+j+1}(t).
\end{equation}
Then $X^2_{c,j}=V_{j-\lceil c/2\rceil}$. For $d\ge2$, on a computable tail,
\begin{align}
 X^d_{c,j}&=X^{d-1}_{c,j}\circ X^d_{c+d},\label{eq:Xrec}\\
 E^d_c&=E^{d-1}_c\circ X^d_{c+d}+E^d_{c+d}.\label{eq:Erec}
\end{align}
\end{lemma}
\begin{proof}
For $d=1$ the first predecessor is finite. For $d=2$, the first
predecessor has $\lfloor(f(t)-c)/2\rfloor$ copies of $\omega$ and a finite
remainder; each $\omega$ copy then costs one replacement. The function $f$
is even, which gives the displayed formula.
For general $d$ write $B=f(t)$ and $B-c=dq+\rho$, $0\le\rho<d$. The first
predecessor has leading part $\omega^{d-1}q$. Its remainder is zero if
$\rho=0$, finite one if $\rho=1$, and $\omega^{\rho-1}$ if $\rho\ge2$.
This is the largest remainder \emph{below $\omega^{d-1}$} under the norm
budget, not a maximum with no ordinal bound. Removing it leaves norm $B$
with spent count one in the first two cases, or norm
$X^{\rho-1}_{B-\rho}(B)$ with count $1+E^{\rho-1}_{B-\rho}(B)$ otherwise.
The remaining $q$ degree-$(d-1)$ blocks have retained masses
$c+d(k-1)$, for $k=q,q-1,\ldots,1$. Their $X$ maps compose and their $E$
counts add. Replacing $c$ by $c+d$ preserves the initial remainder and
initial entry but stops one block earlier. This proves both recurrences,
including the coherent final depth $j$. It is also a rank-recursive exact
evaluator for every displayed kernel and counter.
\end{proof}

\begin{example}[The first counter that must be retained]
For $\lambda=\omega^3$, formula \eqref{eq:prefix} gives
$J_\lambda(t)=X^3_0(t)+E^3_0(t)$. For $\lambda=2\omega^3$, it gives
$J_\lambda(t)=J_{\omega^3}(X^3_4(t))+E^3_4(t)$, since
$N(\omega^3)=4$. Replacing either expression by its outer $X$ alone would
change exact graph equations and could change the result of cancellation.
\end{example}

\subsection{Separation with indices allowed to grow}
We need a uniform estimate, because the block recursion creates retained
masses comparable with $f(t)$. Merely proving a tail for each fixed mass
would not license the induction.

\begin{lemma}[Uniform kernel separation]\label{lem:uniform}
Put $C_d=1000d^2$, $T_3=16$, and for $d\ge4$ put
$T_d=\max(T_{d-1},100C_{d-1}+100d)$. For $d\ge3$, $c,e\ge0$, signed
$j,k$, and $R\ge0$, suppose
\begin{equation}\label{eq:region}
 t\ge T_d,\qquad
 f(t)\ge C_d(c+e+|j|+|k|+R+1).
\end{equation}
Then the depths are defined and at least $f(t)$, and
\begin{align*}
 X^d_{c,j}(t+1)&>f^R(X^d_{e,k}(t)),\\
 c<e\quad&\Longrightarrow\quad
 X^d_{c,j}(t)>f^R(X^d_{e,k}(t)).
\end{align*}
At equal input and equal $c$, exact order is order of $j$. On a common
computed tail, equality of two rank-$d$ values occurs exactly when input,
$c$, and $j$ are identical.
\end{lemma}
\begin{proof}
We first record the elementary degree-two estimate. If
$A=f(w)/2\ge |j|+|k|+R+4$, then
\[
 V_j(w+1)>f^R(V_k(w)).
\]
Indeed $f(w+1)/2-A\ge A$, so the left iterate starts at a larger seed and
has more iterations than the right one plus $R$. At equal input, the $V$
values are ordered exactly by their indices. This argument permits indices
and required depth varying with a later induction parameter, provided the
displayed numerical bound is paid.

For $d=3$, $B=f(t)$ is even and $B\equiv1\pmod3$; write
$a=(B-1)/3$, which is odd. The three initial entries for prefix classes
$\delta=0,1,2$ are
\[
\begin{array}{c|ccc}
\delta&\text{initial }c&\text{entry norm}&\text{spent count}\\ \hline
0&B-1=3a&B&1\\
1&B=3a+1&B&1\\
2&B-2=3(a-1)+2&f(B)&2.
\end{array}
\]
Descending $c$ by three applies $V_{-\lceil(c-3)/2\rceil}$ to the preceding
entry and adds its displayed iteration length to the counter. Write
$Y^\delta_r=X^3_{3r+\delta}(t)$. At the first common entry $r=a-1$,
$Y^0_r=f(Y^1_r)$, since the two subtracted ceilings differ by one. One
further aligned block amplifies this integer input gap to any fixed
$f$-depth by the degree-two estimate. Subsequent indices are favourable:
the class-zero index is no smaller than the class-one index. Next,
$Y^2_{a-1}=f(B)$ whereas
\[
 Y^1_{a-1}=f^\ell(B),\qquad
 \ell=f(B)/2-\lceil(3(a-1)+1)/2\rceil+1.
\]
For $a$ larger than the desired depth by ten, $\ell$ exceeds that depth
plus one. Finally align class two at entry $r-1$ with class zero at entry
$r$. The initial seeds satisfy $f(B)>B$, and
$\lceil(3r-4)/2\rceil\le\lceil(3r-3)/2\rceil$, so the degree-two
estimate again applies in the favourable direction. Consequently
\[
 Y^0_r>f^L(Y^1_r),\quad
 Y^1_r>f^L(Y^2_r),\quad
 Y^2_r>f^L(Y^0_{r+1})
\]
at all required target entries, after paying the fixed depth $L$.
All smaller block entries are at least $B$, and the absolute degree-two
indices are at most $\lceil3a/2\rceil+2$ plus fixed alignment shifts;
$f(B)/2$ pays their sum and $L$. Condition \eqref{eq:region}, with
$C_3=9000$, leaves more than ten blocks beyond the target and the required
$R+|j|+|k|+1$ depth. This proves the same-input statement by chaining the
three adjacent classes.

For cross-input cubic comparison, $B'=f(t+1)>2^B$ supplies more than one
extra leading block. At the smaller process's starting coefficient,
the larger process has completed a quadratic block and has seed at least
$f(B')>f(B)$, exceeding either possible smaller seed. Align the remaining
blocks, shifting the coefficient alignment by at most
$\max(\lfloor c/3\rfloor-\lfloor e/3\rfloor,0)$. Their indices are bounded
by $\lceil3a/2\rceil+2$ plus this shift, and their differences are bounded
by twice the shift plus three. The same degree-two estimate at smaller
entry at least $B$ pays each comparison. Extra remaining blocks only
increase the final norm. This proves the cubic cross-input statement.

For the induction $d\ge4$ set $L=R+|j|+|k|+1$, $B=f(t)$.
Condition \eqref{eq:region} leaves at least $B/(2d)>L+2$ leading
blocks. Every aligned lower-rank mass is at most $2B$, and both its pair
sum and each self-pair sum, together with $L+1$, are at most $5B$.
All smaller aligned entry norms are at least $B$. Since
$B\ge T_d\ge100C_{d-1}$,
\[
 f(B)\ge2^B>5C_{d-1}B.
\]
Thus the induction hypothesis is valid at the \emph{smaller} entry norm
even though the new masses may be comparable with $B$.

At equal $t$ and $c<e$, align the two descents at the initial leading
coefficient of the $e$ process. If no extra leading block was needed,
the $c$ process's remainder is at least as large, with smaller retained
mass. If extra blocks were needed, its last extra full block exceeds the
other initial remainder. These facts order their aligned seeds weakly.
To justify this seed comparison explicitly, maximal predecessors are
monotone in both ordinal bound and norm budget. If entry norms $v\ge w$
and masses $a,b$ satisfy
$f^i(v)-a\ge f^i(w)-b$ at every step, the first process stays above the
second at simultaneous steps, so takes at least as many steps and has at
least as large a terminal norm. Decreasing mass at equal entry is therefore
favourable. A larger entry also pays an adverse mass difference whenever
$f(v)-f(w)\ge a-b$, since successive iterate gaps do not decrease.
At aligned block $k$ the lower masses are $c+d(k-1)<e+d(k-1)$.
Equal seeds use lower-rank same-input separation; distinct seeds use
lower-rank cross-input separation. The estimates just proved pay both,
including lower self-pairs. One such block gives depth $L$, and the rest
preserve it.

At $t+1$, its first budget $B'>2^B$ gives an extra full leading block
before alignment. Its entry is at least $B'$ and retained mass at most
$B'$. The smaller initial remainder is below that block, has entry $B$,
and retained mass at most $B$. The inequality
$f(B')-f(B)>B'$ pays even an adverse retained mass; hence the larger seed
is strict. At subsequent aligned blocks the mass pair is bounded by
$4B$; lower-rank cross-input separation, including its self-pair bounds,
applies as before. This proves the cross-input statement. Monotonicity
at fixed mass extends adjacent comparison to arbitrary larger inputs.

Finally the number of remaining blocks exceeds $|j|+|k|+1$, so signed
depths leave a positive iterate and output at least $f(t)$. Taking the
maximum of the cross-pair and the two self-pair thresholds excludes every
unequal-input or unequal-mass equality. The coherent identity settles the
remaining equal-input, equal-mass case.
\end{proof}

\subsection{Counter leaders, exact ties, and residues}
For $r\ge3$ define $A_r=t$ and successively
\[
 A_{h-1}=X^h_{c+h}(A_h)\quad(h=r,r-1,\ldots,4),\qquad
 u=X^3_{c+3}(A_3),\qquad H=f(u).
\]
We call $H$ the counter's \emph{leader}; $u$ is its leader preimage.

\begin{lemma}[Counter expansion and collision]\label{lem:counter}
On an effective tail,
\begin{equation}\label{eq:counter-expand}
 E^r_c(t)=\frac H2+1-\lceil c/2\rceil
                  +\sum_{h=3}^r E^h_{c+h}(A_h),
 \qquad \frac H2<E^r_c(t)<\frac{3H}{4}.
\end{equation}
Suppose counters of ranks $r\ge s\ge3$ have equal leaders. Then they
have the same prefix index $c$. If $r=s$ their whole values are equal;
if $r>s$ the smaller-rank input is the corresponding $A_s$, and
\begin{equation}\label{eq:collision}
 E^r_c(t)-E^s_c(A_s)=\sum_{h=s+1}^r E^h_{c+h}(A_h).
\end{equation}
Every counter introduced on the right has rank strictly greater than $s$.
For distinct leaders, counter order agrees with leader order. Equal
leaders are ordered by counter rank, with equality at equal rank.
\end{lemma}
\begin{proof}
Iterate \eqref{eq:Erec} until the cubic term; its bottom quadratic count
is $f(X^3_{c+3}(A_3))/2+1-\lceil c/2\rceil$. This proves the exact
identity. A descent with initial norm $w$ and $L$ steps satisfies
$L\le f^L(w)-w$, since each iterate increases its integer argument.
Thus each residual above rank three is less than its next entry norm,
hence less than $u$. The cubic residual plus the fixed correction is also
less than $u$, since it is at most $u-A_3+1-\lceil c/2\rceil$.
It is positive on a computed tail: the number of cubic blocks tends to
infinity. Therefore
\[
 H/2<E^r_c(t)<H/2+(r-2)u<3H/4
\]
once $f(u)>4(r-2)u$ and the finite correction bound is paid.

Equality of leaders first cancels $f$, then compares the two outer cubic
kernels $X^3_{c+3}$. Lemma~\ref{lem:uniform} forces equality of their
prefix indices and arguments. Cancel the equal rank-four kernels next,
then rank five, up to rank $s$. If $r=s$ the original inputs are equal.
If $r>s$ the rank-$s$ input is precisely $A_s$. Subtraction of the exact
recursions gives \eqref{eq:collision}. Its right side is positive and tends
to infinity. For different leaders $H_1<H_2$, their being distinct values
of $f$ gives $H_2\ge2^{H_1+1}$; the half-core bounds make the order strict.
This also proves the assertions with arbitrary fixed outer shifts after
raising the computed tail.
\end{proof}

\begin{lemma}[Effective residues of kernels and counters]\label{lem:kernel-residues}
For each fixed $d,c,j,m$, both $X^d_{c,j}$ and $E^d_c$ have an effectively
computable constant residue modulo $m$ on a computable tail. The $X$
residue is $r_m$. A valid counter residue is
\begin{equation}\label{eq:E-residue}
 E^d_c(t)\equiv
 p_d\bigl((r_{Q_d}-c)\bmod Q_d\bigr)\pmod m,
 \qquad Q_d=d!m^{d-1},
\end{equation}
where $p_d(n)$ counts partitions of $n$ into parts at most $d$.
\end{lemma}
\begin{proof}
A coherent $X$ lies in the range of $f$ with preimage at least $t$ on a
computed tail. Apply Lemma~\ref{lem:residues}.
For the counter, the restricted-partition and congruent-budget theorems
in \cite{wwcounts} prove the following two facts, with full degree
inductions: $Q_d$ is a period of $p_d(n)$ modulo $m$, and a controlled
$\omega^d$ block with retained norm $c$, whose successive budgets are all
congruent to $S$ modulo $Q_d$, has complete count congruent to
$p_d(S-c)$ modulo $m$. Budgets may grow; equality of the budgets is not
assumed. For $t\ge Q_d$, every actual budget $f^{i}(t)$ is congruent to
$r_{Q_d}$, by Lemma~\ref{lem:residues}. Choose $S\ge c$ with that residue
and apply the theorem. The partition period reduces its argument to the
representative in \eqref{eq:E-residue}. Finite dynamic programming computes
$p_d(n)$ for all $n<Q_d$. The domain bound and $t\ge Q_d$ provide an
explicit threshold. This proof of residues uses neither comparison nor
an existential decision procedure.
\end{proof}

\section{A closed existential decision procedure}
Fix a finite maximum degree $D\ge3$. A \emph{word} in one raw variable
is a finite composition of signed constant shifts and the functions
\[
 f,\quad V_j,\quad X^r_{c,j},\quad E^r_c\quad(3\le r\le D),
\]
with fixed integer indices satisfying the preceding conventions. A
word-linear row is an integer linear combination of finitely many such
words, evaluated in possibly different raw variables, plus a constant.
We allow arbitrary Boolean combinations of row comparisons with zero and
fixed-modulus congruences. Each raw variable has an explicit individual
lower bound ensuring that all its words are defined.

Every word has a computable domain tail and a lower offset $M$ with
$W(n)\ge n-M$. This follows structurally: shifts add to the offset;
$f$, $V$, and coherent $X$ dominate their input on their stated tails;
counters do so as well. For the latter, a sufficiently large degree-three
block contains two quadratic blocks, and a quadratic block entered at
least at $f(t)$ already spends more than $f(t)$ steps. Higher-degree blocks
contain one of these. These structural bounds also pay every threshold at
an interior argument of a composition. They are not obtained by treating
the interior argument as a raw variable.

\begin{lemma}[Ordered maximum-variable guards]\label{lem:guards}
For two words $A(x),B(y)$ and the guard $x\ge y$, an effective common
tail permits their three-way comparison to be rewritten as a finite Boolean
combination in which $x$ occurs only linearly and words are applied only
to the remaining variable $y$. For two words at the \emph{same} raw
variable the corresponding tail comparison is effectively decided.
All thresholds and newly generated words are computed before the tail is
used.
\end{lemma}
\begin{proof}
Peel fixed outer shifts. Unequal range values have gaps eventually larger
than a fixed shift. Lemma~\ref{lem:counter} gives the same fact for two
counter values, including equal-leader unequal-rank values whose exact
residual difference tends to infinity. For a counter with preimage $u$
and range value $f(v)$, its half-core bounds give
\[
 E^r_c<f(v)\ \Longleftrightarrow\ u\le v,
 \qquad E^r_c>f(v)\ \Longleftrightarrow\ u>v;
\]
the gaps again tend to infinity. Thus fixed shifts are kept for exact-core
ties and otherwise cannot close the gap.

For $f/f$ strip the two $f$ nodes. For $f$ against a coherent range
kernel, move one fixed depth to its index and strip $f$; for example
$V_j=f\circ V_{j-1}$. Two $V$ heads compare lexicographically by argument,
then by index, using the elementary estimate in the proof of
Lemma~\ref{lem:uniform}. Two equal-rank $X$ heads compare by argument,
then by $(-c,j)$, using that lemma and retaining the exact equal-input
branch.

For $X^r$ against $X^s$ with $r>s\ge3$, expand only the higher-rank head
by \eqref{eq:Xrec} until its outer rank is $s$, and immediately compare
and strip that head against the opposite rank-$s$ head. If $s=2$,
the exposed quadratic head is $V_{j-\lceil c/2\rceil}$ and the same
procedure uses its index comparison. No inverse counter is introduced.
For counter/range comparison use the displayed preimage rule. For
counter/counter comparison use their leader preimages followed by the
rank and exact-tie rule of Lemma~\ref{lem:counter}.

Here is a termination measure, needed because mixed-rank expansion can
increase the written expression. First count all $E$ nodes in the two
words. Replacing a counter by its leader-preimage composition removes
its outer $E$ node and retains its argument once. For the range part,
assign each rank-$s$ occurrence weight $\omega^{D-s}$, where $f$ has
rank one and $V$ rank two, and naturally sum these weights. A mixed
comparison deletes the opposite rank-$s$ occurrence and introduces only
finitely many occurrences of ranks strictly above $s$, retaining the
original higher-rank occurrence. Its ordinal weight decreases. Equal-rank
stripping and $f$ cancellation decrease weight as well. The lexicographic
pair (number of $E$ nodes, range weight) is therefore well-founded along
every recursive child. There are finitely many branches at each step.

Ultimately the $x$ side is a shifted leaf, so its comparison with a word
at $y$ is linear in $x$; or the $y$ side is a shifted leaf and the $x$
side is nonlinear. In the latter case the nonlinear value is at least
$2^{x-M}$ for a computed $M$, so eventually exceeds $y$ plus any fixed
shift, since $y\le x$. When both words have the same raw variable,
continue stripping until leaves or leaf/nonlinear growth decides the
comparison; do not create a purported linear atom $x<W(x)$.
Take the finite maximum of all local bounds using the structural offsets
of interior arguments. The finite recursion tree, rather than a guessed
fixed envelope of indices, computes that maximum.
\end{proof}

The guards compare individual word values. A linear row can still cancel
its largest values, so it requires a different induction.

\begin{lemma}[Ranked cancellation]\label{lem:row}
For a finite collection of word-linear rows and a chosen maximal raw
variable $x$, one can compute a tail and finite guarded cases in which
each row either has a known strict sign or is exactly linear in $x$.
The guards can themselves be made linear in $x$ by
Lemma~\ref{lem:guards}. No correlation between word values is discarded.
\end{lemma}
\begin{proof}
Absorb fixed outer shifts into the constant and keep raw-linear terms
separate. The other top-level atoms are ordinary range values (rank zero
for this row argument) and counters of ranks $3,\ldots,D$. Partition their
ordinary values and counter leaders into ordered equality classes, retaining
every equality and order guard. Merge equal-rank counters in a class as
\emph{whole counters}, justified by Lemma~\ref{lem:counter}. Prefer a
representative actually involving $x$ when one remains in the current
class. A class now has ordinary coefficient $a$ and at most one coefficient
$b_r$ for each counter rank. Its leading coefficient is
\[
 a+\tfrac12\sum_r b_r.
\]
If it vanishes, pay its residuals before deciding any sign. When $a\ne0$,
expand the counters by \eqref{eq:counter-expand}; the common half-leader
cancels the ordinary atom, leaving only residual counters and constants.
When $a=0$, let $s$ be the least surviving counter rank. Rewrite each
higher counter by \eqref{eq:collision} as the common rank-$s$ counter plus
higher-rank residuals. Since $\sum b_r=0$, the common term cancels. Remove
an entirely zero class. Repartition after every rewrite, since its newly
exposed leaders may coincide with leaders elsewhere in the row.

Count \emph{top-level} row atoms, assigning the ordinary rank-zero atom
weight $\omega^D$ and a rank-$r$ counter weight $\omega^{D-r}$.
Natural sum is a well-founded measure. The first rewrite loses a rank-zero
atom and adds only positive-rank atoms. The second loses the least rank
$s$ and adds only ranks strictly greater than $s$. Existing higher-rank
atoms can proliferate, but cannot offset that ordinal decrease. Nested
nodes inside their arguments are not included in this measure; their
comparisons terminate by the separate guard measure. Every branching row
rewrite therefore strictly decreases weight. The tree is finitely
branching and has no infinite branch, so it is finite. Enumerating it to
completion computes all generated words, coefficients, prefix indices,
and guard thresholds before any finite-prefix case is split off.

After the rewrites every surviving leading coefficient has absolute value
at least $1/2$. Let $P=f(u)$ be the greatest surviving leader, $S$ the
current absolute nonlinear coefficient mass, and $A,B$ bounds for the
absolute raw-linear coefficients and constant. The exact counter
corrections have total size at most $DSu$. Every smaller leader $Q$
satisfies $Q<\log_2 P$, so its total contribution is at most
$S\log_2P$. For $u\ge8DS+16$ each of these two errors is less than $P/16$.

If a current nonlinear $x$ occurrence survives, the greatest leader is
at least its leader. Structural offsets therefore give $u\ge x-M$ for a
computed $M$. Require $x\ge2M+16$ and $x\ge8(A+B+1)$. The elementary
inequality $2^{x-M}\ge x^2$ on this range, proved at the initial value
and then by doubling, gives $Ax+B<P/8$. Total error is less than $P/4$,
while the leading magnitude is at least $P/2$. Hence its coefficient
strictly determines the sign. If no current nonlinear $x$ occurrence
remains, keep the exact row, which is linear in $x$.

Growth provenance refers only to current occurrences. A deleted leader
passes no $x$ label to a residual which does not actually contain $x$.
This prevents an invalid sign decision after cancellation has removed all
nonlinear dependence on $x$. Take maxima of these explicit bounds over the
finite row and guard trees. Ordered partitions retain the original shared
arguments throughout, even when an inconsistent branch is later discarded.
\end{proof}

\begin{theorem}[Closed existential word arithmetic]\label{thm:existential}
The existential closure of every Boolean combination of word-linear row
comparisons and fixed-modulus congruences in the finite-degree language
above is decidable, on its explicit word domains.
\end{theorem}
\begin{proof}
First remove a variable which occurs only linearly, if any, by ordinary
Presburger quantifier elimination. Words in other variables are parameters
in this step; the eliminated variable occurs in none of them. The output
is still a Boolean combination of the same kind. Repeat or evaluate the
closed matrix if no variables remain.

Otherwise branch on a maximal raw variable $x$, retaining $x\ge x_i$.
Compute the complete finite guard and row trees of
Lemmas~\ref{lem:guards}--\ref{lem:row}. Compute residue thresholds for each
congruence as well: Lemma~\ref{lem:kernel-residues}, composed with interior
lower-offset bounds, makes every nonlinear word's residue constant.
Collect shifts in its fixed residue. On that tail, congruences are thus
congruences in raw variables only. Take a common $N$ paying all original
and generated domains, sign estimates, guards, and residue bounds.

Before executing any tail transformation, split the \emph{original current
matrix} into the all-raw-inputs-$\ge N$ branch and finitely many branches
fixing some raw $x_i=c<N$ above its original lower bound. In a finite
branch substitute $c$ into its original words and evaluate them exactly.
Do not evaluate a newly introduced coherent negative-depth word below the
tail on which it was introduced. Exact kernel evaluation is finite by the
block recursion. Each such branch removes a quantified variable.

On the all-large branch, use the guarded row and congruence transformations.
Everything now contains $x$ only linearly. Presburger elimination of
$\exists x$, treating words in the remaining variables as parameters,
again produces word-linear rows and congruences with one fewer variable.
Preserve all Boolean connectives and every equality/order guard. No word
has been applied to a sum, and no auxiliary derived value has been made
an independent witness.

The induction measure for the whole algorithm is the number of raw
existential variables. It strictly decreases in every finite-prefix,
pure-linear, and all-large branch. The subsidiary guard and row trees are
finite by their own measures. Consequently this is a terminating effective
procedure. No feasible complexity estimate follows: finite exception values
and arithmetic thresholds can be extremely large.
\end{proof}

This theorem is deliberately about \emph{closed existential} sentences.
In a formula with a free parameter, that parameter could be the largest raw
argument and could not be removed by the stated induction. We have proved
neither free-parameter quantifier elimination nor the full first-order
theory of arbitrary polynomial-collapse functions.

\section{Exact polynomial graphs and the internal-level-two consumer}
\begin{theorem}[Polynomial collapse graphs]\label{thm:graphs}
For each fixed nonzero polynomial limit $\lambda<\omega^\omega$, the graph
$y=h_\lambda(n)$ is effectively expressible as a finite disjunction of
word-linear equations in $n,y$, with explicit domains and exact finite
exceptions. One finite degree bound suffices for each finite collection of
such graphs. Arbitrary fixed nesting and addition of these functions admit
a total functional existential graph translation.
\end{theorem}
\begin{proof}
Remove one unit at the least positive degree of $\lambda$, writing
$\lambda=\mu+\omega^d$, and put $c=N(\mu)$. There are finitely many
monomial units. Starting with $2J_0(t)=2t$, compile \eqref{eq:prefix}
recursively: substitute the unary word $X^d_c(t)$ for the input in the
already compiled expression for $2J_\mu$, and add $2E^d_c(t)$. For $d=1$
use $X=f$, $2E=2$; for $d=2$ use
\[
 X=V_{-\lceil c/2\rceil},\qquad
 2E=f(t)-2\lceil c/2\rceil+2;
\]
for $d\ge3$ retain the $X$ and $E$ heads. This substitution composes unary
words only. It never applies a function to a sum of compiled values.

Finally replace $t$ by $n+N(\lambda)$. On a computed tail this yields
$2y=P_\lambda(n)$ with an integer word-linear $P_\lambda$. The finitely
many excluded $n\ge0$ are evaluated by the total closed predecessor
algorithm and included as disjuncts $n=k$, $y=h_\lambda(k)$. The graph is
therefore total and functional on the original natural domain.
To flatten a nested expression, introduce a fresh natural variable for
each arithmetic node: addition and fixed multiples have linear equations;
a collapse node has the graph just constructed, with its original argument
variable. Preserve all repeated original variable names and quantify every
new node variable. This creates no independence assumption: the graph
equations enforce precisely the original correlations.
\end{proof}

\begin{theorem}[Canonical source comparison through level two]\label{thm:leveltwo}
Universal weak, WW-strict, and equality comparison is decidable for original
canonical $T$ assignments with ground free variables and maximum internal
type level at most two. The maximum includes every subterm and binder type,
not merely the result and free-variable types.
\end{theorem}
\begin{proof}
The actual-source induction in \cite{wwfoundations} gives the forward
residual translation. Its level bound can be seen directly in that generated
grammar. Coordinate two has no exponent constructor: one would require a
coordinate three. Its closed values are therefore natural linear
combinations of $1$ and $\omega$. Coordinate one forms natural sums and
products and powers
$2^{a\omega+b}=\omega^a2^b$ for finite $a,b$; its values are polynomials
in $\omega$. Each direct scalar head $\omega\otimes c$ is consequently
a polynomial limit below $\omega^\omega$. Temporary functional binder
variables disappear under the canonical abstraction substitution, and the
internal-level hypothesis includes their types. No statement about arbitrary
ambient coordinate classes is needed.

Evaluate and compare closed positive coordinates. Translate a zeroth
counterexample to the requested universal relation using
Theorem~\ref{thm:graphs}; all fresh graph variables are quantified, with
nonnegative domains, and original variables remain positive. The result
is a closed existential Boolean combination in a finite-degree word
language. Theorem~\ref{thm:existential} decides whether such a counterexample
exists. This separately decides each of the source's strict, weak, and
equality relations with its actual coordinate convention. There is no
fixed degree bound across all input terms: the finite degree needed is
computed from each pair.
\end{proof}

\section{Arbitrary fixed heads: sparse endpoints}
Polynomial graphs were needed for completely unrestricted arrangements of
polynomial heads. Different arrangements permit arbitrary fixed heads
without constructing their graphs. This section isolates the spacing fact
used by all of them.

\begin{lemma}[Sparse endpoint comparison]\label{lem:endpoints}
For a fixed nonzero limit $\lambda$, $J_\lambda$ is strictly increasing.
Its distinct norm inputs have distinct terminal endpoints, ordered in the
same way. For any two distinct endpoints $Z<W$ one has
\begin{equation}\label{eq:sparse}
 W\ge2^{Z+1}.
\end{equation}
Every value with endpoint $Z$ lies strictly between $Z$ and $2Z$ and below
the next value of $f$.
\end{lemma}
\begin{proof}
Strict increase is the controlled finite-tail strictness proved in
\cite{wwfoundations}: when $t<u$, the smaller ordinal has norm
$t\le f(u)$, so the controlled comparison premise holds. Write
$J_\lambda(t)=f^L(t)+L$ and $J_\lambda(u)=f^M(u)+M$.
If their endpoints were equal, injectivity of $f$ would force $L>M$;
otherwise cancelling $L$ iterates would give $t\ge u$. The equal endpoints
and $L>M$ would imply $J_\lambda(t)>J_\lambda(u)$, a contradiction.
Put $s=f^{L-1}(t)$. Since $t\ge N(\lambda)\ge2$, one has $L\le s-1$,
so $f(s)<J_\lambda(t)<2f(s)<f(s+1)$ by Lemma~\ref{lem:growth}.
These intervals are disjoint and preserve order. If $W>Z=f(s)$ is another
endpoint, then $W\ge f(s+1)\ge2^{Z+1}$.
\end{proof}

\subsection{Nested common-head sums}
Fix any finite list of fixed nonzero limits $\lambda_1,\ldots,\lambda_l$.
Let $\mathcal A_0$ consist of natural constants, positive raw variables,
addition, and $f$, with arbitrary fixed nesting. Recursively set
\begin{equation}\label{eq:common-class}
 \mathcal A_r=\left\{c+\sum_i a_i h_{\lambda_r}(B_i):
 c\ge0,\ a_i>0,\ B_i\in\mathcal A_{r-1}\right\}.
\end{equation}
The sum can be empty. Each input supplies this finite level certificate.
There is no open remainder outside the head sum at a positive level.
For example $x+h_\lambda(y)$ does not pass that positive-level certificate
merely because both $x,y$ are allowed at level zero. Heads at different
levels may be completely different, and may exceed $\omega^\omega$.

\begin{theorem}[Common-head layers]\label{thm:common}
For $E,F\in\mathcal A_r$, each pointwise relation $E<F$, $E=F$, $E>F$
has an effectively constructed formula in $(\N;0,1,+,<,f)$ with the same
raw free variables and exactly the same truth at every positive valuation.
Hence arbitrary Boolean and quantified combinations of certified comparisons
are decidable.
\end{theorem}
\begin{proof}
Subtract the two expressions and write the row as
\[
 Q=K+\sum_{i<m}a_i J_{\lambda_r}(t_i),\qquad
 t_i=N(\lambda_r)+B_i.
\]
Coefficients and $K$ are signed integers. With no remaining atoms the sign
is the computable sign of $K$. Otherwise set
$C=\sum_i|a_i|$ and $L=\max(4,2C,4|K|+4)$.
On the branch with all $t_i\ge L$, partition the inputs by equal values,
combine coefficients, and delete zero blocks. If every block cancels,
$Q=K$. Otherwise its sign is the coefficient of the greatest surviving
input. Indeed let $W$ be its endpoint and $Z$ the greatest lower endpoint,
if any. Lemma~\ref{lem:endpoints} gives
$W\ge2^{Z+1}\ge2Z^2\ge4CZ$, since $Z\ge L\ge\max(4,2C)$.
All lower contributions have total absolute value less than
$2CZ\le W/2$. Also $|K|<W/4$, whereas the leading contribution has
absolute value greater than $W$. Its sign cannot change. Larger canceled
blocks play no part in this estimate. Enumerating the finitely many total
preorders of the labeled inputs produces a finite formula, guarded by
$t_i=t_j$, $t_i<t_j$, and $t_i\ge L$.

For the other branch choose a small remaining index, enumerate
$N(\lambda_r)\le c<L$, retain its exact guard $t_i=c$, evaluate the
closed $J_{\lambda_r}(c)$, fold $a_iJ_{\lambda_r}(c)$ into $K$, and
remove that occurrence. Recompute $L$ and repeat. Each small call removes
one occurrence; its possibly enormous updated constant is still a fixed
computable integer. Overlapping branches are harmless in a disjunction.
No open value is folded into a constant before its whole argument is fixed.

Induct now on $r$. At level zero comparisons are base-language atoms.
At higher levels every guard compares preceding-level expressions with
each other or with constants; a proposed constant below $N(\lambda_r)$
gives a constant truth value. Thus every generated predicate is covered
by the induction hypothesis. The lexicographic pair (level, remaining
head occurrences) proves termination: small folding decreases the second,
guard translation the first. Shared variables remain shared throughout.
The result is an exact base formula, to which
Theorem~\ref{thm:snyatkov} applies after any desired raw quantifiers are
added.
\end{proof}

\begin{example}[Cancellation changes the controlling input]
Consider $2h_\lambda(x)+h_\lambda(y)-3h_\lambda(z)$. On the computed
large-input branch, a unique largest argument determines the sign by its
coefficient. But on the guard $x=z$, the row becomes exactly
$h_\lambda(y)-h_\lambda(x)$, whose sign is that of $y-x$; when all three
arguments agree the row is zero. The algorithm retains those guards rather
than treating three head values as independent large numbers. Small inputs
are evaluated by the finite folding procedure, so this is an exact
comparison theorem rather than only an eventual rule.
\end{example}

\subsection{Different heads in one control interval}
\begin{lemma}[Head order in a strict interval]\label{lem:head-window}
Let $\lambda>\mu$ be nonzero fixed limits. If
$t\ge\max(N(\lambda),N(\mu)+2)$ and
$N(\mu)\le u<f(t)$, then $J_\lambda(t)>J_\mu(u)$ and their endpoints
differ. Consequently, if finitely many norm inputs $t_i$ have minimum
$s\ge\max_iN(\lambda_i)+2$ and $\max_i t_i<f(s)$, their outputs are
ordered lexicographically by (ordinal head, norm input). Equality requires
equality of both, and different output blocks have distinct sparse endpoints.
\end{lemma}
\begin{proof}
The least limit above $\mu$ is $\mu+\omega$, of norm $N(\mu)+2$.
Controlled strictness at common norm $t$, or equality of the two heads,
gives
\[
 J_\lambda(t)\ge J_{\mu+\omega}(t)
       =J_\mu(f(t))+1>J_\mu(u).
\]
The exact middle identity follows by removing the final $\omega$ block:
its first predecessor retains $\mu$, assigns the remaining full budget
to its finite tail, and has total norm $f(t)$. If the two displayed
outputs had equal endpoints, their complete counts would satisfy
$L>M$; canceling the common $f$ iterates would give
$u=f^{L-M}(t)\ge f(t)$, a contradiction. For the finite collection, the
lower input is below $f(s)\le f(t)$ for every higher-head input $t$.
Apply this argument to different heads and
Lemma~\ref{lem:endpoints} within the same head.
\end{proof}

The interval inequality is strict for a reason. At $u=f(t)$,
\[
 J_{2\omega}(t)=f^2(t)+2,\qquad
 J_\omega(u)=f^2(t)+1.
\]
Their endpoints coincide, and
$J_{2\omega}(t)-2J_\omega(u)=-f^2(t)<0$ despite the positive coefficient
of the greatest ordinal head. An algorithm using head order at that
boundary would therefore be incorrect.

\section{Balanced heterogeneous layers}
The next certificate permits different sibling heads and an open lower-level
remainder. At level zero use the same $\mathcal A_0$ as before. At a
positive level choose a finite common preceding-level basis
$B_1,\ldots,B_k$. It may have repeated or dependent entries. New expressions
are nonnegative affine combinations of the $B_j$ and atoms
\begin{equation}\label{eq:balanced-atoms}
 h_\lambda\!\left(c+\sum_j a_jB_j\right),
       \qquad c\ge0,\quad a_j\ge1.
\end{equation}
Each atom may have its own arbitrary fixed head and coefficients. The
essential condition is that \emph{every head argument includes every
common basis entry positively}. The basis is checked on the whole pair
being compared, including its outside remainder. A variable appearing
only in that remainder is not admitted by ignoring it while choosing the
basis. Zero multiples are discarded when computing actual supports.

\begin{theorem}[Balanced comparison]\label{thm:balanced}
For two expressions with one finite balanced certificate, their three-way
pointwise comparison has an effectively constructed base formula in
$(\N;0,1,+,<,f)$ with exactly the same free variables. Arbitrary Boolean
and quantified combinations of such comparisons are decidable.
\end{theorem}
\begin{proof}
A signed comparison row has the form
\begin{equation}\label{eq:balanced-row}
 Q=K+\sum_jr_jB_j+\sum_{i<m}q_iJ_{\lambda_i}(t_i),\qquad
 t_i=b_i+\sum_j a_{ij}B_j,\quad b_i=N(\lambda_i)+c_i.
\end{equation}
The head-free case is already a signed affine comparison at the preceding
level. Otherwise put
\[
 S=\sum_jB_j,\quad C=\sum_i|q_i|,\quad C_0=\max_j|r_j|,
 \quad A=\max_{i,j}a_{ij},\quad B=\max_i b_i,
\]
with $C_0=0$ for an absent remainder and $A=1$ for an empty basis. Set
\begin{equation}\label{eq:balanced-L}
 L=\max\{16,2C,4C_0+4|K|+4,A+B+2,\max_iN(\lambda_i)+2\}.
\end{equation}
On the branch $s=\min_i t_i\ge L$ one has
\[
 S\le t_i\le AS+B,\qquad
 \max_i t_i\le As+B<s^2\le2^s\le f(s).
\]
Thus Lemma~\ref{lem:head-window} orders all head outputs by head and then
input. Partition the inputs within each equal head, combine coefficients,
and discard zero blocks. If all cancel, the remaining row is exactly the
preceding-level $K+\sum r_jB_j$.

If a block survives, let $W$ be the endpoint of the last one and $v$ its
norm input. Its coefficient has absolute value at least one, so its
contribution has magnitude greater than $W$. The sparse-endpoint estimate
bounds all lower head contributions by $W/2$, exactly as in the proof of
Theorem~\ref{thm:common}. The outside remainder is bounded using this
\emph{surviving} input, not a possibly canceled maximum:
\[
 \left|K+\sum_jr_jB_j\right|\le |K|+C_0v
          <v^2/4\le f(v)/4\le W/4.
\]
The surviving last coefficient therefore determines the strict sign.
The finitely many guarded preorders within equal heads express this rule
by comparisons of affine preceding-level expressions.

We must retain every small exception. Every positive expression in the
certificate satisfies $B(\vec x)\ge x$ for each raw variable in its
actual support. Induct through addition and $f$, and at a head use
$h_\lambda(e)>e$: the finite ordinal $e$ is below $\lambda+e$ and has
norm $e\le f(N(\lambda)+e)$, so controlled strictness and $\psi(e)=e$
apply. Fixed positive multiples preserve the property. If some $t_i<L$,
positivity of every $a_{ij}$ implies $S<L$. Therefore every raw variable
anywhere in the basis, including the outside remainder, is less than $L$.
Enumerate the finite box $[1,L-1]^V$ on its \emph{actual} support $V$;
evaluate each closed row and its inputs exactly, retaining precisely
points where some $t_i<L$ and the requested sign holds. Variables outside
$V$ remain unrestricted; empty support gives a single empty valuation.

All large-branch guards, and the all-canceled remainder, are affine
comparisons of the $B_j$ and thus have their preceding-level certificate.
The finite-point exception formula is already in the base language.
Induction on certificate level proves the claimed translation and
termination. Theorem~\ref{thm:snyatkov} supplies its arithmetic decision.
Repeated basis entries and shared variables never become independent
parameters in this construction.
\end{proof}

\begin{example}[A lower-head argument may be larger]
Let $E,F$ have a common preceding-level certificate and $\lambda>\mu$.
The row
\[
 3h_\lambda(E+F+2)-2h_\mu(2E+F+1)-F
\]
is balanced on basis $(E,F)$. Its two norm arguments need not have the
same order as their ordinal heads, and $F$ remains open outside them.
Nevertheless, once both norms exceed \eqref{eq:balanced-L}, the row is
positive by the higher-head coefficient. Its exact comparison is completed
by the finite support box. Omitting $F$ from one head argument would
remove the inference that a small input bounds the entire open remainder.
\end{example}

\section{One arbitrary core: direct Presburger comparison}
An especially useful specialization no longer requires a certificate for the
interior core. Let $E$ be any expression made from natural constants,
positive raw variables, addition, fixed nonnegative multiples, and arbitrary
fixed nonzero-limit functions $h_\lambda$, with arbitrary finite nesting.
No subtraction or exposed $f$ constructor occurs inside this stated grammar.
The only property needed below is $E\ge x$ for every variable in its actual
support, proved by the positive-source induction just given. It need not
have a balanced-level certificate, nor range independently over all naturals.

\begin{theorem}[Arbitrary single-core rows]\label{thm:single-core}
For the signed row
\begin{equation}\label{eq:single-core}
 Q(\vec x)=K+rE(\vec x)+
       \sum_i q_i h_{\lambda_i}(a_iE(\vec x)+c_i),
 \quad a_i\ge1,\ c_i\ge0,
\end{equation}
with fixed integer $K,r,q_i$ and fixed nonzero limit heads, each sign
predicate $Q<0$, $Q=0$, $Q>0$ has an effectively computable finite
Presburger formula in the original raw variables. Thus universal comparison
of these rows needs no arithmetic decision theorem about $f$.
\end{theorem}
\begin{proof}
Collect identical triples $(\lambda_i,a_i,c_i)$ and discard their zero
coefficients. If any heads remain, use \eqref{eq:balanced-L} with the
single formal basis $E$, $C_0=|r|$, and $b_i=N(\lambda_i)+c_i$. Increase
that bound to
\[
 T=\max\bigl(L,1+\max_{i,j}|b_i-b_j|\bigr).
\]
For $E\ge T$, the norms lie in one strict $f$ interval. Within equal
heads their affine inputs are ordered lexicographically by $(a_i,b_i)$:
a difference of integer slopes dominates the fixed offset difference.
Equal triples have already been collected. Lemma~\ref{lem:head-window}
and the remainder estimate in the balanced proof therefore show that
the sign is the coefficient of the greatest triple
$(\lambda_i,a_i,b_i)$. This use of the balanced \emph{numerical estimate}
requires no certificate or comparison algorithm for $E$ itself.
If no heads remain, use $T=\max(1,|K|+1)$; the final sign of $rE+K$ is
$\sgn(r)$ when $r\ne0$, otherwise $\sgn(K)$. Denote the computed final
sign by $\sigma$.

Let $V$ be the actual raw support of $E$. Empty support is handled by one
closed evaluation. Otherwise, outside $[1,T-1]^V$ there is an $x\ge T$,
and positivity gives $E\ge T$. Evaluate every row in the finite box exactly
and record its sign. For a desired sign $s$ the exact formula is
\[
 \left(\bigvee_{x\in V}x\ge T\right)\wedge[\sigma=s]
 \quad\vee\quad
 \bigvee_{\substack{v\in[1,T-1]^V\\\sgn Q(v)=s}}
       \bigwedge_{x\in V}x=v(x).
\]
It is a finite Presburger formula on the original variables. The box is
effective and each closed evaluation terminates. If the tail falsifies a
universal comparison, a concrete admissible countervaluation sets one
supported variable to $T$ and all others to one; the dominance property
certifies it without evaluating a giant collapse value.
\end{proof}

For example, even a core $E=h_\alpha(x)+h_\beta(h_\gamma(y)+z)$ outside
a proposed balanced certificate may be used in
$3h_\lambda(E+c)-2h_\mu(2E+d)-E$. Its support is still $(x,y,z)$.
This result is a comparison theorem for that row, not a graph definition
or a general entry-map algorithm for the core. An unrelated open term
outside $E$, a variable ordinal head, or a nonlinear dependence on $E$
is outside the statement.

\section{What these decisions do and do not establish}
All source consumers begin by computing the literal canonical assignment.
They then compare its closed positive coordinates and check the appropriate
scalar certificate. This direction is essential: none asserts that every
residual expression in its mathematical grammar is the image of a $T$ term.
A supplied certificate may make a recognition procedure incomplete as a
search for all possible decompositions, but its check and the ensuing
comparison are exact on accepted inputs.

The finite-coefficient theorem includes all recursor-free ground-open
sources at arbitrary internal type levels. The polynomial theorem includes
all ground-open sources whose \emph{entire} internal syntax has type level
at most two, including recursors. At arbitrary internal levels, common-head,
balanced, or single-core certificates can admit some nonpolynomial heads.
Every result keeps repeated variables and correlated subexpressions, exact
prefix masses, and all complete descent counters.

The mathematical algorithms are total; they are not promises of practical
execution. A resource-limited implementation which cannot evaluate a closed
exception or complete the inherited arithmetic procedure must report that
unmet computation. A symbolic name for a closed value is sufficient where
it cancels exactly, but is not a numerical evaluation when its sign is
needed. We have neither implemented Snyatkov's full procedure here nor
computed any literal giant $F_5$ value.

The remaining ground-open region includes unbalanced multivariate
arrangements of arbitrary nonpolynomial heads, without any supplied
certificate bringing them into the proved classes. The whole ground-open
weak and strict comparison problem remains open. The type-confined ordering
and non-ground expressivity companions address other parts and boundaries
of the original effectiveness question \cite{wwordering,wwexpressivity};
one-variable comparison has its own stronger methods
\cite{wwunivariate}. None of these restricted results, individually or
combined, announces a full resolution of the Wilken--Weiermann question.

\begin{thebibliography}{99}

\bibitem{ww}Gunnar Wilken and Andreas Weiermann.
Derivation lengths classification of G\"odel's T extending Howard's assignment.
\emph{Logical Methods in Computer Science} \textbf{8}(1:19) (2012), 1--44.
\href{https://doi.org/10.2168/LMCS-8(1:19)2012}
{doi:10.2168/LMCS-8(1:19)2012}.
\href{https://lmcs.episciences.org/1073}{Official article and full text}.
The 2012 revised article is the version used throughout.

\bibitem{snyatkov}A.~S. Snyatkov.
Decidability of the theory
$T_f=\operatorname{Th}(\omega,0,1,<,+,f(x),F(x))$.
\emph{Vestnik Tverskogo Gosudarstvennogo Universiteta,
Series Applied Mathematics}, no.~2(9), pp.~39--51, 2008.
In Russian.
\url{https://eprints.tversu.ru/id/eprint/494/1/19950136_2008_2_snyatkov.pdf}.

\bibitem{wwfoundations}\emph{The Wilken--Weiermann Comparison Problem: Foundations and a Guide to the Results}. Companion working manuscript, 2026.\
\href{../05-wilken-weiermann-foundations/paper.pdf}{Included manuscript}.

\bibitem{wwordering}\emph{A Type-Confined Termination Ordering for G\"odel's T}. Companion working manuscript, 2026.\
\href{../06-godel-t-termination-ordering/paper.pdf}{Included manuscript}.

\bibitem{wwunivariate}\emph{Equality and One-Variable Comparison for Positive Collapse Expressions}. Companion working manuscript, 2026.\
\href{../07-collapse-equality-and-one-variable/paper.pdf}{Included manuscript}.

\bibitem{wwexpressivity}\emph{Actual-Source Arithmetic and Undecidability Boundaries for Canonical Comparison}. Companion working manuscript, 2026.\
\href{../09-canonical-source-expressivity/paper.pdf}{Included manuscript}.

\bibitem{wwcounts}\emph{Controlled Ordinal Descents, Restricted Partitions and Cartier Observers}. Companion working manuscript, 2026.\
\href{../10-collapse-descent-and-modular-counts/paper.pdf}{Included manuscript}.
\end{thebibliography}

\end{document}
